\section{Разработка информационно-поисковой системы}
\label{sec:development}

\subsection{Выбор средств реализации}

Система реализована по клиент-серверной схеме и разворачивается в локальной вычислительной сети средствами Docker Compose. При разработке использован следующий технологический стек.

\begin{itemize}
	\item \textbf{Python 3.11} "--- язык серверной части: предобработка текста, вероятностная модель, REST API, оценка качества;
	\item \textbf{FastAPI} "--- веб-фреймворк; маршрутизация HTTP-запросов, валидация входных данных моделями Pydantic, автоматическая документация OpenAPI;
	\item \textbf{Uvicorn} "--- ASGI-сервер для запуска FastAPI;
	\item \textbf{PostgreSQL 16} "--- реляционная СУБД; хранит документы, инвертированный индекс, эталонную разметку и журналы;
	\item \textbf{psycopg2} "--- драйвер PostgreSQL для Python;
	\item \textbf{pymorphy3} "--- морфологический анализатор русского языка; лемматизация словоформ;
	\item \textbf{matplotlib} "--- построение графиков метрик качества;
	\item \textbf{pypdf}, \textbf{python-docx} "--- извлечение текста из PDF- и DOCX-файлов;
	\item \textbf{React 18}, \textbf{Vite}, \textbf{react-router-dom} "--- одностраничный веб-интерфейс;
	\item \textbf{nginx} "--- раздача собранного фронтенда и проксирование запросов \texttt{/api} на бэкенд;
	\item \textbf{Docker Compose} "--- контейнеризация трёх сервисов и воспроизводимое развёртывание в ЛВС.
\end{itemize}

Вероятностная модель выбрана в соответствии с вариантом задания. Расчёт RSV по формуле~\eqref{eq:bm25} выполняется SQL-запросом на стороне СУБД: индекс уже лежит в таблицах \texttt{terms} и \texttt{postings}, и нет необходимости выгружать коллекцию в память приложения.

FastAPI выбран как лёгкий фреймворк с явной схемой запросов: тело поискового запроса, параметры модели и флажок обратной связи проверяются до входа в бизнес-логику. React обеспечивает интерактивный интерфейс "--- элемент искусственного интеллекта по варианту~5: автодополнение, исправление опечаток и подсветка найденных слов не требуют перезагрузки страницы.

Контейнеризация закрывает требование о сфере применения <<локальная вычислительная сеть>>: сервис публикуется на порту 8080 и доступен всем узлам сети по адресу узла, на котором запущен Docker.

\subsection{Описание используемых готовых компонентов}

\begin{itemize}
	\item \textbf{pymorphy3} "--- для каждого русскоязычного токена выбирается наиболее вероятная нормальная форма. Словарь \texttt{pymorphy3-dicts-ru} встраивается в Docker-образ, сетевые обращения при работе не требуются;
	\item \textbf{FastAPI / Uvicorn} "--- REST API и раздача описания эндпоинтов. CORS открыт, поскольку фронтенд и бэкенд в разработке могут работать на разных портах;
	\item \textbf{PostgreSQL} "--- расчёт BM25 выполняется одним запросом с \texttt{unnest} массивов идентификаторов и весов терминов. Агрегат \texttt{array\_agg} возвращает список совпавших лемм для подсветки;
	\item \textbf{matplotlib} "--- бэкенд \texttt{Agg} строит PNG без графической подсистемы. Для интерфейса графики кодируются в base64, для отчёта сохраняются в каталог \texttt{png/};
	\item \textbf{pypdf / python-docx} "--- подключаются только при загрузке соответствующего формата;
	\item \textbf{React} "--- клиентские страницы поиска, документов, загрузки, оценки качества и справки. Взаимодействие с сервером "--- \texttt{fetch} по JSON API;
	\item \textbf{nginx} "--- \texttt{try\_files} отдаёт SPA, \texttt{proxy\_pass} направляет \texttt{/api/} на сервис \texttt{backend};
	\item \textbf{Docker Compose} "--- оркестрация сервисов \texttt{db}, \texttt{backend} и \texttt{frontend} с проверкой готовности базы данных.
\end{itemize}

\subsection{Структура разработанной системы}

Система состоит из трёх уровней: клиентского (React SPA за nginx), серверного (FastAPI) и уровня данных (PostgreSQL). Общая схема приведена на рисунке~\ref{fig:architecture}.

\begin{figure}[H]
	\centering
	\includegraphics[width=1.0\textwidth]{architecture.png}
	\caption{Структура информационно-поисковой системы.}
	\label{fig:architecture}
\end{figure}

Серверная часть разделена на модули по ответственности:

\begin{itemize}
	\item \texttt{main.py} "--- точка входа, REST-эндпоинты и жизненный цикл приложения (миграции, загрузка корпуса);
	\item \texttt{text\_processing.py} "--- токенизация, лемматизация, фильтрация стоп-слов;
	\item \texttt{indexer.py} "--- построение и обновление инвертированного индекса;
	\item \texttt{document\_loader.py} "--- извлечение текста из файлов разных форматов;
	\item \texttt{probabilistic.py} "--- веса RSJ, BM25, псевдорелевантная обратная связь, baseline TF-IDF;
	\item \texttt{query\_assistant.py} "--- автодополнение, исправление опечаток, расширение синонимами, сниппеты;
	\item \texttt{search\_service.py} "--- оркестрация поиска, отсечение хвоста, журналирование;
	\item \texttt{evaluator.py} "--- метрики качества и графики;
	\item \texttt{migrate.py} "--- применение SQL-миграций;
	\item \texttt{bootstrap.py} "--- первичная загрузка коллекции и эталонной разметки.
\end{itemize}

Клиентский уровень содержит пять страниц: поиск, список документов, загрузка, оценка качества и справка. Основной сценарий "--- обработка поискового запроса "--- показан на рисунке~\ref{fig:seq-search}.

\begin{figure}[H]
	\centering
	\includegraphics[width=1.0\textwidth]{sequence_search.png}
	\caption{Диаграмма последовательности обработки поискового запроса.}
	\label{fig:seq-search}
\end{figure}

Развёртывание выполняется командой \texttt{docker compose up -{}-build} из каталога программной части. При старте контейнер бэкенда дожидается готовности PostgreSQL, применяет миграции и, если база пуста, индексирует тестовую коллекцию. Интерфейс доступен по адресу \texttt{http://<host>:8080}.

\subsection{Структура базы данных системы}

Схема содержит восемь таблиц, перечисленных в таблице~\ref{tab:db-tables}. Диаграмма сущностей приведена на рисунке~\ref{fig:db}.

\begin{table}[H]
	\centering
	\begin{tabular}{|p{3.4cm}|p{8.6cm}|}
		\hline
		\textbf{Таблица} & \textbf{Назначение} \\
		\hline
		\texttt{documents} & Документы коллекции: заголовок, текст, саммари, путь к файлу, длина в леммах, дата добавления \\
		\hline
		\texttt{terms} & Словарь: нормальная форма термина и документная частота \\
		\hline
		\texttt{postings} & Инвертированный индекс: частота термина в документе \\
		\hline
		\texttt{search\_logs} & Журнал запросов: исходный текст, нормализованный образ, число результатов, время \\
		\hline
		\texttt{eval\_queries} & Эталонные запросы для оценки качества \\
		\hline
		\texttt{eval\_qrels} & Разметка релевантности (qrels) по имени файла документа \\
		\hline
		\texttt{eval\_runs} & Сохранённые сводки прогонов оценки \\
		\hline
		\texttt{schema\_migrations} & Учёт применённых SQL-миграций \\
		\hline
	\end{tabular}
	\caption{Таблицы базы данных и их назначение.}
	\label{tab:db-tables}
\end{table}

\begin{figure}[H]
	\centering
	\includegraphics[width=1.0\textwidth]{db_schema.png}
	\caption{Схема базы данных системы.}
	\label{fig:db}
\end{figure}

Связь эталонной разметки с документами выполняется по полю \texttt{source\_file}, а не по числовому идентификатору: при пересоздании базы идентификаторы могут измениться, имена файлов корпуса "--- нет.

\subsection{Основные алгоритмы реализации компонентов}

\textbf{Предобработка.} Функция \texttt{normalize} возвращает список значимых лемм (листинг~\ref{lst:normalize}). Позиции токенов сохраняются, чтобы сниппет мог подсветить исходные словоформы, а не только нормальные формы.

\begin{lstlisting}[caption={Конвейер предобработки русскоязычного текста.},label={lst:normalize}]
def normalize(text: str) -> list[str]:
    return [t.lemma for t in tokenize(text) if is_meaningful(t.lemma)]

def is_meaningful(lemma: str) -> bool:
    return len(lemma) >= MIN_TOKEN_LEN and lemma not in STOPWORDS
\end{lstlisting}

\textbf{Индексирование.} При добавлении документа частоты лемм записываются в \texttt{postings}, словарь пополняется новыми терминами, затем пересчитываются документные частоты (рисунок~\ref{fig:seq-index}). Повторная загрузка файла с тем же \texttt{source\_file} обновляет документ, а не создаёт дубликат.

\begin{figure}[H]
	\centering
	\includegraphics[width=1.0\textwidth]{sequence_indexing.png}
	\caption{Диаграмма последовательности индексирования документа.}
	\label{fig:seq-index}
\end{figure}

Поддерживаемые форматы загрузки перечислены в таблице~\ref{tab:formats}.

\begin{table}[H]
	\centering
	\begin{tabular}{|p{4.6cm}|p{7.4cm}|}
		\hline
		\textbf{Расширение} & \textbf{Механизм извлечения текста} \\
		\hline
		\texttt{.txt}, \texttt{.md}, \texttt{.csv}, \texttt{.log} & декодирование UTF-8, CP1251, KOI8-R, Latin-1 \\
		\hline
		\texttt{.pdf} & библиотека pypdf \\
		\hline
		\texttt{.docx} & библиотека python-docx \\
		\hline
		\texttt{.html}, \texttt{.htm} & удаление тегов, script- и style-блоков \\
		\hline
		\texttt{.rtf} & удаление управляющих слов RTF \\
		\hline
	\end{tabular}
	\caption{Поддерживаемые форматы загружаемых документов.}
	\label{tab:formats}
\end{table}

\textbf{Веса и BM25.} Вес без обратной связи и расчёт RSV на стороне СУБД приведены в листингах~\ref{lst:rsj} и~\ref{lst:bm25}.

\begin{lstlisting}[caption={Сглаженный вес Робертсона -- Спарк Джонс.},label={lst:rsj}]
def rsj_weight(n_docs: float, df: int) -> float:
    return math.log(1.0 + (n_docs - df + 0.5) / (df + 0.5))
\end{lstlisting}

\begin{lstlisting}[caption={Расчёт RSV по BM25 в PostgreSQL.},label={lst:bm25}]
SELECT p.doc_id,
       SUM(
           qt.w * (p.tf * (k1 + 1.0))
           / (p.tf + k1 * (1.0 - b + b * d.doc_len / stats.avgdl))
       ) AS score
FROM postings p
JOIN qt ON qt.term_id = p.term_id
JOIN documents d ON d.id = p.doc_id
CROSS JOIN stats
GROUP BY p.doc_id
ORDER BY score DESC
\end{lstlisting}

\textbf{Интеллектуальный интерфейс.} По варианту~5 основным реализуемым элементом является интерфейс с пользователем. Его функции:

\begin{itemize}
	\item автодополнение по префиксу из словаря коллекции (запрос \texttt{/api/suggest});
	\item исправление опечаток расстоянием Дамерау -- Левенштейна не более двух операций;
	\item расширение запроса синонимами и аббревиатурами из тезауруса (<<ИИ>> $\to$ <<искусственный интеллект>>);
	\item сниппет с HTML-подсветкой найденных слов и отдельный список словоформ запроса, встретившихся в документе;
	\item фильтр выдачи по дате добавления документа;
	\item встроенная справка на отдельной странице.
\end{itemize}

Последовательность работы подсказок и исправления опечаток показана на рисунке~\ref{fig:seq-assist}.

\begin{figure}[H]
	\centering
	\includegraphics[width=1.0\textwidth]{sequence_assistant.png}
	\caption{Диаграмма последовательности интеллектуального интерфейса запроса.}
	\label{fig:seq-assist}
\end{figure}

\begin{lstlisting}[caption={Расстояние Дамерау -- Левенштейна.},label={lst:damerau}]
current[j] = min(
    current[j - 1] + 1,      # insert
    previous[j] + 1,         # delete
    previous[j - 1] + cost,  # replace
)
if i > 1 and j > 1 and a[i - 1] == b[j - 2] and a[i - 2] == b[j - 1]:
    current[j] = min(current[j], before_previous[j - 2] + cost)
\end{lstlisting}

Страница поиска отображает исправленный запрос (<<Возможно, вы имели в виду...>>), добавленные синонимы, веса терминов и для каждого документа "--- активную ссылку, сниппет и список найденных слов. Это закрывает требование задания о составе поисковой выдачи.

\subsection{Информация о тестовой коллекции документов}

Тестовая коллекция состоит из 18 русскоязычных текстов объёмом 300"---600 слов, относящихся к пяти тематическим группам: естественные науки, история, спорт, техника, информационные технологии. Разнородность подобрана сознательно: на узкой однотемной коллекции вероятностная и векторная модели плохо различимы, потому что документы делят большую часть словаря. Перечень документов приведён в таблице~\ref{tab:collection}.

\begin{table}[H]
	\centering
	\begin{tabular}{|c|p{3.6cm}|p{8.0cm}|}
		\hline
		\textbf{ID} & \textbf{Файл} & \textbf{Тематика} \\
		\hline
		1 & aviaciya & гражданская авиация, подъёмная сила крыла \\
		\hline
		2 & bazy\_dannyh & реляционные базы данных и SQL \\
		\hline
		3 & chernye\_dyry & чёрные дыры, гравитация, гравитационные волны \\
		\hline
		4 & drevniy\_rim & Древний Рим и римское право \\
		\hline
		5 & elektromobili & электромобили и тяговые аккумуляторы \\
		\hline
		6 & futbol & футбол: правила и история \\
		\hline
		7 & genetika\_dnk & генетика, ДНК, редактирование генома \\
		\hline
		8 & hokkey & хоккей с шайбой \\
		\hline
		9 & iskusstvennyy\_intellekt & искусственный интеллект и его применение \\
		\hline
		10 & kosmos\_issledovaniya & исследование космоса, спутники, Марс \\
		\hline
		11 & lokalnye\_seti & локальные вычислительные сети \\
		\hline
		12 & neyronnye\_seti & нейронные сети и машинное обучение \\
		\hline
		13 & olimpiyskie\_igry & Олимпийские игры \\
		\hline
		14 & operacionnye\_sistemy & операционные системы и ядро \\
		\hline
		15 & robototehnika & промышленная робототехника \\
		\hline
		16 & vakciny\_immunitet & вакцинация и иммунная система \\
		\hline
		17 & velikoe\_knyazhestvo & Великое княжество Литовское \\
		\hline
		18 & vtoraya\_mirovaya & Вторая мировая война \\
		\hline
	\end{tabular}
	\caption{Тестовая коллекция документов.}
	\label{tab:collection}
\end{table}

Документы загружаются автоматически при пустой базе из каталога \texttt{backend/app/corpus/} в алфавитном порядке имён файлов: идентификатор в таблице~\ref{tab:collection} совпадает с полем \texttt{id} в базе. Эталонная разметка хранится в файле \texttt{qrels.json}: 18 запросов с указанием релевантных файлов. После индексации разметка привязывается к числовым идентификаторам через \texttt{source\_file}.

\subsection{Результаты тестирования системы}

Тестирование проводилось автоматически сценарием \texttt{backend/scenarios.py} и вручную через веб-интерфейс. Проверялись следующие сценарии:

\begin{enumerate}
	\item предобработка: лемматизация словоформ, отсев стоп-слов, сохранение латиницы;
	\item инициализация: 18 документов, словарь из 2091 термина, 3313 постингов, отсутствие дубликатов при повторном старте;
	\item вероятностные веса: редкий термин получает больший вес, чем частый (2,539 против 0,204 на контрольной паре);
	\item поиск: однословный запрос, многословный запрос, запрос из слов, отсутствующих в коллекции;
	\item интерфейс: исправление опечатки <<нейроные>> $\to$ <<нейронный>>, расширение аббревиатуры <<ИИ>>, автодополнение по префиксу, подсветка сниппета, список найденных словоформ;
	\item загрузка: извлечение текста из HTML, отказ неподдерживаемого формата с понятным сообщением;
	\item индексирование и удаление: новый документ становится доступен поиску и исчезает из индекса после удаления;
	\item оценка качества: расчёт таблиц и графиков по эталонным запросам.
\end{enumerate}

Все 20 автоматических проверок завершились успешно. Средняя длина документа составила 232,9 леммы.

\subsection{Результаты оценки по метрикам}

Оценка выполнялась на 18 эталонных запросах. Для каждого запроса система возвращала выдачу с отсечением хвоста, после чего считались метрики~\eqref{eq:prf}, $P@k$, MAP и nDCG@10. Последовательность расчёта показана на рисунке~\ref{fig:seq-metrics}. Результаты по запросам для модели BM25 без обратной связи приведены в таблице~\ref{tab:metrics}.

\begin{figure}[H]
	\centering
	\includegraphics[width=1.0\textwidth]{sequence_metrics.png}
	\caption{Диаграмма последовательности оценки качества поиска.}
	\label{fig:seq-metrics}
\end{figure}

\begin{table}[H]
	\centering
	\small
	\begin{tabular}{|p{5.6cm}|c|c|c|c|c|}
		\hline
		\textbf{Запрос} & \textbf{P} & \textbf{R} & \textbf{F1} & \textbf{P@5} & \textbf{AP} \\
		\hline
		исследование космоса и планет & 1,000 & 1,000 & 1,000 & 0,400 & 1,000 \\
		\hline
		чёрные дыры и гравитационные волны & 1,000 & 1,000 & 1,000 & 0,200 & 1,000 \\
		\hline
		вакцинация и иммунитет человека & 1,000 & 1,000 & 1,000 & 0,200 & 1,000 \\
		\hline
		наследственность гены и днк & 1,000 & 1,000 & 1,000 & 0,200 & 1,000 \\
		\hline
		вторая мировая война & 0,333 & 1,000 & 0,500 & 0,200 & 1,000 \\
		\hline
		древний рим и римское право & 0,500 & 1,000 & 0,667 & 0,200 & 1,000 \\
		\hline
		великое княжество литовское & 1,000 & 1,000 & 1,000 & 0,200 & 1,000 \\
		\hline
		олимпийские игры & 0,333 & 1,000 & 0,500 & 0,200 & 1,000 \\
		\hline
		футбол правила игры & 0,250 & 1,000 & 0,400 & 0,200 & 1,000 \\
		\hline
		хоккей с шайбой & 0,500 & 1,000 & 0,667 & 0,200 & 1,000 \\
		\hline
		электромобиль и тяговый аккумулятор & 1,000 & 1,000 & 1,000 & 0,200 & 1,000 \\
		\hline
		подъёмная сила крыла самолёта & 1,000 & 1,000 & 1,000 & 0,200 & 1,000 \\
		\hline
		промышленный робот манипулятор & 1,000 & 1,000 & 1,000 & 0,200 & 1,000 \\
		\hline
		нейронные сети и машинное обучение & 0,500 & 1,000 & 0,667 & 0,400 & 1,000 \\
		\hline
		реляционная база данных и sql & 0,500 & 1,000 & 0,667 & 0,200 & 1,000 \\
		\hline
		локальная вычислительная сеть & 0,333 & 1,000 & 0,500 & 0,200 & 1,000 \\
		\hline
		операционная система и ядро & 1,000 & 1,000 & 1,000 & 0,200 & 1,000 \\
		\hline
		спортивные соревнования & 1,000 & 0,667 & 0,800 & 0,400 & 0,667 \\
		\hline
		\textbf{Среднее} & \textbf{0,736} & \textbf{0,981} & \textbf{0,798} & \textbf{0,233} & \textbf{0,981} \\
		\hline
	\end{tabular}
	\caption{Метрики качества BM25 по тестовым запросам.}
	\label{tab:metrics}
\end{table}

Сравнение трёх моделей ранжирования приведено в таблице~\ref{tab:models}.

\begin{table}[H]
	\centering
	\begin{tabular}{|l|c|c|c|c|c|c|}
		\hline
		\textbf{Модель} & \textbf{P} & \textbf{R} & \textbf{F1} & \textbf{P@5} & \textbf{MAP} & \textbf{nDCG@10} \\
		\hline
		BM25 & 0,736 & 0,981 & 0,798 & 0,233 & 0,981 & 0,987 \\
		\hline
		BM25 + PRF & 0,584 & 0,981 & 0,665 & 0,233 & 0,981 & 0,987 \\
		\hline
		TF-IDF & 0,594 & 0,981 & 0,680 & 0,233 & 0,981 & 0,987 \\
		\hline
	\end{tabular}
	\caption{Сравнение моделей ранжирования на тестовой коллекции.}
	\label{tab:models}
\end{table}

Графическое представление "--- столбчатая диаграмма точности, полноты и F1 по запросам (рисунок~\ref{fig:metrics-q}), 11-точечная кривая полноты-точности (рисунок~\ref{fig:pr}) и сравнение моделей (рисунок~\ref{fig:models}).

\begin{figure}[H]
	\centering
	\includegraphics[width=1.0\textwidth]{metrics_by_query.png}
	\caption{Точность, полнота и F1-мера по эталонным запросам (BM25).}
	\label{fig:metrics-q}
\end{figure}

\begin{figure}[H]
	\centering
	\includegraphics[width=0.92\textwidth]{pr_curve.png}
	\caption{11-точечная интерполированная кривая полноты-точности.}
	\label{fig:pr}
\end{figure}

\begin{figure}[H]
	\centering
	\includegraphics[width=1.0\textwidth]{model_comparison.png}
	\caption{Сравнение BM25, BM25 с обратной связью и TF-IDF.}
	\label{fig:models}
\end{figure}

\subsection{Анализ результатов и предложения по улучшению}

Анализ полученных данных показывает следующее.

\begin{itemize}
	\item Полнота близка к единице (0,981): на коллекции из 18 документов релевантные тексты почти всегда попадают в выдачу. Исключение "--- широкий запрос <<спортивные соревнования>>, для которого не найден один из трёх эталонных документов.
	\item Точность ограничивается общими словами запроса. Для <<футбол правила игры>> слово <<игра>> встречается в нескольких спортивных текстах, и отсечение хвоста оставляет четыре документа при одном эталонном ($P = 0{,}250$). Специфичные запросы с редкими терминами (<<чёрные дыры>>, <<тяговый аккумулятор>>) дают $P = 1$.
	\item MAP и nDCG@10 почти равны единице: когда релевантный документ найден, он стоит на первой позиции. Качество ранжирования на данной коллекции выше, чем качество отсечения нерелевантного хвоста.
	\item $P@5$ и $P@10$ занижены по построению: в знаменателе всегда стоит $k$, тогда как релевантных документов обычно один-два, а выдача после отсечения короче пяти позиций.
	\item BM25 превосходит TF-IDF по точности (0,736 против 0,594) и F1-мере (0,798 против 0,680). Насыщение частоты и поправка на длину уменьшают вклад случайных совпадений.
	\item Псевдорелевантная обратная связь на этой коллекции вредна: точность падает до 0,584. При $R = 5$ в псевдорелевантное множество попадают документы соседних тем, и пересчёт весов усиливает общие слова.
\end{itemize}

Предложения по улучшению работы системы:

\begin{enumerate}
	\item расширить коллекцию до сотен документов, чтобы полнота и $P@k$ стали информативнее, а обратная связь "--- полезнее;
	\item учитывать биграммы и устойчивые словосочетания (<<вторая мировая война>>, <<локальная вычислительная сеть>>);
	\item подбирать порог отсечения и параметр $R$ по валидационным запросам, а не фиксировать их заранее;
	\item заменить ручной тезаурус векторными эмбеддингами или готовым словарём синонимов;
	\item добавить оценку по кликам пользователя как истинную обратную связь вместо псевдорелевантной;
	\item вынести расчёт метрик в регулярный прогон при изменении корпуса.
\end{enumerate}

\subsection{Вывод по разделу}

Реализована информационно-поисковая система с вероятностной моделью Okapi BM25, инвертированным индексом в PostgreSQL и интеллектуальным веб-интерфейсом. Система поддерживает загрузку документов, поиск по естественно-языковому запросу с подсветкой найденных слов, автодополнение, исправление опечаток и программную оценку качества. На тестовой коллекции средняя F1-мера BM25 составила 0,798 при полноте 0,981, что выше результата базовой векторной модели.
